% Hidden-state capture (split from the old 3-panel methods figure).
\begin{tikzpicture}[
  font=\footnotesize,
  >={Stealth[length=2mm]},
  node distance=4mm and 5mm,
  databox/.style = {draw=orange!75!black, fill=orange!8, rounded corners=2pt,
                    inner sep=4pt, align=center},
  opbox/.style   = {draw=blue!60!black,   fill=blue!6,   rounded corners=2pt,
                    inner sep=4pt, align=center},
  outbox/.style  = {draw=green!50!black,  fill=green!8,  rounded corners=2pt,
                    inner sep=4pt, align=center},
  basebox/.style = {draw=violet!60!black, fill=violet!8, rounded corners=2pt,
                    inner sep=4pt, align=center},
  arr/.style     = {->, semithick, gray!50!black},
]

\node [databox, text width=3.4cm] (prompt) at (0,0)
       {Prompt\\\textit{``Who won the Nobel\\Prize in Physics in 1921?''}};

\node [opbox, right=6mm of prompt, text width=3.6cm] (lm)
      {Subject LM\\(e.g.\ Qwen3-1.7B; $L{=}28$)\\
       \footnotesize\itshape prefill: one forward pass,\\nothing generated yet};

\node [outbox, right=6mm of lm, text width=3.8cm] (htensor)
      {$h^{(\ell)}_T$ at every layer $\ell$\\$(n,\ L{+}1,\ d)$\\$=(784,\ 29,\ 2048)$};

\draw [arr] (prompt) -- (lm);
\draw [arr] (lm)     -- (htensor);

\end{tikzpicture}
